\documentclass[10pt,a4paper]{amsart}
\usepackage{amsmath,amssymb,amsthm,graphicx}
\usepackage{verbatim}
\newenvironment{reference}{\comment}{\endcomment}
%\newenvironment{reference}{}{}
\newenvironment{slogan}{\comment}{\endcomment}
\newenvironment{history}{\comment}{\endcomment}

% For commutative diagrams we use Xy-pic
\usepackage[all]{xy}

% We use 2cell for 2-commutative diagrams.
\xyoption{2cell}
% Theorem environments.
%
\theoremstyle{plain}
\newtheorem{theorem}[subsection]{Theorem}
\newtheorem{proposition}[subsection]{Proposition}
\newtheorem{lemma}[subsection]{Lemma}

\theoremstyle{definition}
\newtheorem{definition}[subsection]{Definition}
\newtheorem{example}[subsection]{Example}
\newtheorem{exercise}[subsection]{Exercise}
\newtheorem{situation}[subsection]{Situation}

\theoremstyle{remark}
\newtheorem{remark}[subsection]{Remark}
\newtheorem{remarks}[subsection]{Remarks}

\numberwithin{equation}{subsection}

% Macros
%
\def\lim{\mathop{\mathrm{lim}}\nolimits}
\def\colim{\mathop{\mathrm{colim}}\nolimits}
\def\Spec{\mathop{\mathrm{Spec}}}
\def\Hom{\mathop{\mathrm{Hom}}\nolimits}
\def\Ext{\mathop{\mathrm{Ext}}\nolimits}
\def\SheafHom{\mathop{\mathcal{H}\!\mathit{om}}\nolimits}
\def\SheafExt{\mathop{\mathcal{E}\!\mathit{xt}}\nolimits}
\def\Sch{\mathit{Sch}}
\def\Mor{\mathop{\mathrm{Mor}}\nolimits}
\def\Ob{\mathop{\mathrm{Ob}}\nolimits}
\def\Sh{\mathop{\mathit{Sh}}\nolimits}
\def\NL{\mathop{N\!L}\nolimits}
\def\CH{\mathop{\mathrm{CH}}\nolimits}
\def\proetale{{pro\text{-}\acute{e}tale}}
\def\etale{{\acute{e}tale}}
\def\QCoh{\mathit{QCoh}}
\def\Ker{\mathop{\mathrm{Ker}}}
\def\Im{\mathop{\mathrm{Im}}}
\def\Coker{\mathop{\mathrm{Coker}}}
\def\Coim{\mathop{\mathrm{Coim}}}

% Boxtimes
%
\DeclareMathSymbol{\boxtimes}{\mathbin}{AMSa}{"02}

%
% Macros for moduli stacks/spaces
%
\def\QCohstack{\mathcal{QC}\!\mathit{oh}}
\def\Cohstack{\mathcal{C}\!\mathit{oh}}
\def\Spacesstack{\mathcal{S}\!\mathit{paces}}
\def\Quotfunctor{\mathrm{Quot}}
\def\Hilbfunctor{\mathrm{Hilb}}
\def\Curvesstack{\mathcal{C}\!\mathit{urves}}
\def\Polarizedstack{\mathcal{P}\!\mathit{olarized}}
\def\Complexesstack{\mathcal{C}\!\mathit{omplexes}}
% \Pic is the operator that assigns to X its picard group, usage \Pic(X)
% \Picardstack_{X/B} denotes the Picard stack of X over B
% \Picardfunctor_{X/B} denotes the Picard functor of X over B
\def\Pic{\mathop{\mathrm{Pic}}\nolimits}
\def\Picardstack{\mathcal{P}\!\mathit{ic}}
\def\Picardfunctor{\mathrm{Pic}}
\def\Deformationcategory{\mathcal{D}\!\mathit{ef}}
\usepackage{cleveref}
\crefname{originallemma}{original lemma}{original lemmas}
\Crefname{originallemma}{Original lemma}{Original lemmas}
\makeatletter
\newlabel{algebra-lemma-minimal-over-1}{{00KV}{source}{}{sourcepointer-1}{}}
\newlabel{algebra-lemma-minimal-over-1@cref}{{[originallemma][0][]00KV}{[1][0][]source}{}{}{}}
\makeatother
\crefname{originallemma}{original lemma}{original lemmas}
\Crefname{originallemma}{Original lemma}{Original lemmas}
\makeatletter
\newlabel{decent-spaces-lemma-codimension-local-ring}{{0ED1}{source}{}{sourcepointer-2}{}}
\newlabel{decent-spaces-lemma-codimension-local-ring@cref}{{[originallemma][0][]0ED1}{[1][0][]source}{}{}{}}
\makeatother
\crefname{originallemma}{original lemma}{original lemmas}
\Crefname{originallemma}{Original lemma}{Original lemmas}
\makeatletter
\newlabel{decent-spaces-lemma-decent-generic-points}{{0ABV}{source}{}{sourcepointer-3}{}}
\newlabel{decent-spaces-lemma-decent-generic-points@cref}{{[originallemma][0][]0ABV}{[1][0][]source}{}{}{}}
\makeatother
\crefname{originalreference}{original reference}{original references}
\Crefname{originalreference}{Original reference}{Original references}
\makeatletter
\newlabel{situation-setup}{{02QL}{source}{}{sourcepointer-4}{}}
\newlabel{situation-setup@cref}{{[originalreference][0][]02QL}{[1][0][]source}{}{}{}}
\makeatother
\crefname{originallemma}{original lemma}{original lemmas}
\Crefname{originallemma}{Original lemma}{Original lemmas}
\makeatletter
\newlabel{spaces-divisors-lemma-characterize-effective-Cartier-divisor}{{083C}{source}{}{sourcepointer-5}{}}
\newlabel{spaces-divisors-lemma-characterize-effective-Cartier-divisor@cref}{{[originallemma][0][]083C}{[1][0][]source}{}{}{}}
\makeatother
\crefname{originaldefinition}{original definition}{original definitions}
\Crefname{originaldefinition}{Original definition}{Original definitions}
\makeatletter
\newlabel{spaces-over-fields-definition-Weil-divisor}{{0ENJ}{source}{}{sourcepointer-6}{}}
\newlabel{spaces-over-fields-definition-Weil-divisor@cref}{{[originaldefinition][0][]0ENJ}{[1][0][]source}{}{}{}}
\makeatother
\crefname{originaldefinition}{original definition}{original definitions}
\Crefname{originaldefinition}{Original definition}{Original definitions}
\makeatletter
\newlabel{topology-definition-dimension-function}{{02I9}{source}{}{sourcepointer-7}{}}
\newlabel{topology-definition-dimension-function@cref}{{[originaldefinition][0][]02I9}{[1][0][]source}{}{}{}}
\makeatother
\title{Prime divisors and components of effective Cartier divisors have delta-dimension one less than the ambient space}
\author{The Stacks Project Authors}
\date{stacks-project:spaces-chow.tex:local-snapshot-20261012; GNU FDL 1.2 or later}
\begin{document}
\maketitle
\noindent\textit{Editorial note.} All delivered mathematical text is the author's own, copied from the pinned Stacks Project chapter file spaces-chow.tex; no statement or proof step was rewritten, reordered or removed.\par\medskip
\noindent\textit{Editorial note.} The counted proof is the author's complete argument: an unpacking of the definitions for the first part and, for the second, an etale neighbourhood cutting the divisor by a nonzerodivisor together with the principal ideal theorem\par\medskip
\noindent\textit{Editorial note.} The label lemma-divisor-delta-dimension names the counted statement, and the proof environment immediately following it in the source is delivered verbatim as its proof.\par\medskip
\noindent\textit{Editorial note.} The counted statement is carried by the author's own lemma environment in the chapter source.\par\medskip
\noindent\textit{Allowed standard prerequisites (editorial).}\begin{itemize}
\item The definition of a prime divisor on an integral good algebraic space and of the associated Weil divisor, and the identification of codimension one points with the local rings of dimension one
\item The definition of a dimension function and its compatibility with codimension, so that codimension one means delta-dimension one less
\item The local description of an effective Cartier divisor as the zero scheme of a nonzerodivisor on an etale neighbourhood, and Krull's principal ideal theorem
\end{itemize}\medskip
\ifdefined\hypertarget\hypertarget{sourcepointer-1}{}\fi
\noindent\textit{Source reference.} Original lemma 00KV: sources/stacks-chow/tags.txt, source lines 765--765. Prerequisite pointer invoked at proof line(s) 1627; it names a standard fact of a sibling chapter that the counted argument uses, and every step of the counted proof built on it is retained verbatim in the delivered body.\par\medskip
\ifdefined\hypertarget\hypertarget{sourcepointer-2}{}\fi
\noindent\textit{Source reference.} Original lemma 0ED1: sources/stacks-chow/tags.txt, source lines 17976--17976. Prerequisite pointer invoked at proof line(s) 1613; it names a standard fact of a sibling chapter that the counted argument uses, and every step of the counted proof built on it is retained verbatim in the delivered body.\par\medskip
\ifdefined\hypertarget\hypertarget{sourcepointer-3}{}\fi
\noindent\textit{Source reference.} Original lemma 0ABV: sources/stacks-chow/tags.txt, source lines 13033--13033. Prerequisite pointer invoked at proof line(s) 1625; it names a standard fact of a sibling chapter that the counted argument uses, and every step of the counted proof built on it is retained verbatim in the delivered body.\par\medskip
\ifdefined\hypertarget\hypertarget{sourcepointer-4}{}\fi
\noindent\textit{Source reference.} Original reference 02QL: sources/stacks-spaces-chow/pinned-original-source.tex, source lines 264--264. Prerequisite pointer invoked at proof line(s) 1594; it names a standard fact of the same chapter that the counted argument uses, and every step of the counted proof built on it is retained verbatim in the delivered body.\par\medskip
\ifdefined\hypertarget\hypertarget{sourcepointer-5}{}\fi
\noindent\textit{Source reference.} Original lemma 083C: sources/stacks-chow/tags.txt, source lines 10282--10282. Prerequisite pointer invoked at proof line(s) 1623; it names a standard fact of a sibling chapter that the counted argument uses, and every step of the counted proof built on it is retained verbatim in the delivered body.\par\medskip
\ifdefined\hypertarget\hypertarget{sourcepointer-6}{}\fi
\noindent\textit{Source reference.} Original definition 0ENJ: sources/stacks-chow/tags.txt, source lines 18344--18344. Prerequisite pointer invoked at proof line(s) 1612; it names a standard fact of a sibling chapter that the counted argument uses, and every step of the counted proof built on it is retained verbatim in the delivered body.\par\medskip
\ifdefined\hypertarget\hypertarget{sourcepointer-7}{}\fi
\noindent\textit{Source reference.} Original definition 02I9: sources/stacks-chow/tags.txt, source lines 3292--3292. Prerequisite pointer invoked at proof line(s) 1615; it names a standard fact of a sibling chapter that the counted argument uses, and every step of the counted proof built on it is retained verbatim in the delivered body.\par\medskip
\begin{lemma}
\label{lemma-divisor-delta-dimension}
In Situation \ref{situation-setup} let $X/B$ be good. Assume $X$ is
integral.
\begin{enumerate}
\item If $Z \subset X$ is an integral closed subspace, then
the following are equivalent:
\begin{enumerate}
\item $Z$ is a prime divisor,
\item $|Z|$ has codimension $1$ in $|X|$, and
\item $\dim_\delta(Z) = \dim_\delta(X) - 1$.
\end{enumerate}
\item If $Z$ is an irreducible component of an effective Cartier
divisor on $X$, then $\dim_\delta(Z) = \dim_\delta(X) - 1$.
\end{enumerate}
\end{lemma}

\begin{proof}
Part (1) follows from the definition of a prime divisor
(Spaces over Fields, Definition
\ref{spaces-over-fields-definition-Weil-divisor}),
Decent Spaces, Lemma \ref{decent-spaces-lemma-codimension-local-ring},
and the definition of a dimension function
(Topology, Definition \ref{topology-definition-dimension-function}).

\medskip\noindent
Let $D \subset X$ be an effective Cartier divisor. Let $Z \subset D$
be an irreducible component and let $\xi \in |Z|$ be the generic point.
Choose an \'etale neighbourhood $(U, u) \to (X, \xi)$ where $U = \Spec(A)$
and $D \times_X U$ is cut out by a nonzerodivisor $f \in A$, see
Divisors on Spaces, Lemma
\ref{spaces-divisors-lemma-characterize-effective-Cartier-divisor}.
Then $u$ is a generic point of $V(f)$ by
Decent Spaces, Lemma \ref{decent-spaces-lemma-decent-generic-points}.
Hence $\mathcal{O}_{U, u}$ has dimension $1$ by
Krull's Hauptidealsatz (Algebra, Lemma \ref{algebra-lemma-minimal-over-1}).
Thus $\xi$ is a codimension $1$ point on $X$ and $Z$ is a prime
divisor as desired.
\end{proof}

\end{document}
